% =============================================================================
% Theory Development Draft: Bounded Latent Transport Framework
% Status: WORKING DRAFT - Not for submission
% =============================================================================
\documentclass[letterpaper]{article}
\usepackage{amsmath,amssymb,amsthm}
\usepackage{bm,bbm}
\usepackage{algorithm,algpseudocode}
\usepackage[colorlinks,linkcolor=blue,citecolor=blue]{hyperref}
\usepackage[margin=1in]{geometry}
\usepackage{enumitem}

% ---- theorem environments ----
\newtheorem{proposition}{Proposition}[section]
\newtheorem{lemma}[proposition]{Lemma}
\newtheorem{corollary}[proposition]{Corollary}
\newtheorem{definition}[proposition]{Definition}
\newtheorem{assumption}[proposition]{Assumption}
\newtheorem{remark}[proposition]{Remark}
\newtheorem{example}[proposition]{Example}

% ---- math shorthand ----
\newcommand{\R}{\mathbb{R}}
\newcommand{\E}{\mathbb{E}}
\newcommand{\Var}{\operatorname{Var}}
\newcommand{\SWD}{\operatorname{SWD}}
\newcommand{\OT}{\operatorname{OT}}
\newcommand{\cA}{\mathcal{A}}
\newcommand{\cL}{\mathcal{L}}
\newcommand{\cC}{\mathcal{C}}
\newcommand{\cH}{\mathcal{H}}
\newcommand{\cN}{\mathcal{N}}
\newcommand{\cP}{\mathcal{P}}
\newcommand{\cB}{\mathcal{B}}
\newcommand{\cE}{\mathcal{E}}
\newcommand{\id}{\operatorname{id}}
\newcommand{\tr}{\operatorname{tr}}
\newcommand{\diag}{\operatorname{diag}}
\newcommand{\KL}{\operatorname{KL}}
\newcommand{\TV}{\operatorname{TV}}
\newcommand{\argmin}{\operatorname{argmin}}
\newcommand{\argmax}{\operatorname{argmax}}
\newcommand{\supp}{\operatorname{supp}}
\newcommand{\sort}{\operatorname{sort}}
\newcommand{\proj}{\operatorname{proj}}

\title{Theory Development Draft:\\A Bounded Latent Transport Framework\\for OT-Coupled Flow Matching}
\author{Internal Working Document}
\date{\today}

\begin{document}
\maketitle

\begin{abstract}
This document develops a bounded formal analysis of the OT-coupled latent flow matching framework.
We provide explicit propositions with assumptions, proofs, and direct correspondence to implementation.
The results are organized as a coherent theory package spanning problem formalization (Proposition~0),
conditional velocity estimation (Proposition~1), path-energy interpretation (Proposition~2),
discretization error analysis (Proposition~3), endpoint distribution control (Proposition~4),
and OT-coupling stabilization (Proposition~5).

\textbf{Disclaimer.}
These results do not claim exact equivalence to a full Schr\"odinger Bridge or Benamou--Brenier solver.
They provide a bounded formal justification for the three-part design of endpoint construction,
path learning, and endpoint correction, grounded in the actual implementation.
\end{abstract}

\tableofcontents

\newpage

%%%%%%%%%%%%%%%%%%%%%%%%%%%%%%%%%%%%%%%%%%%%%%%%%%%%%%%%%%%%%%%%%%%%%%%%%%%%%%%%
\section{Notation and Preliminaries}
%%%%%%%%%%%%%%%%%%%%%%%%%%%%%%%%%%%%%%%%%%%%%%%%%%%%%%%%%%%%%%%%%%%%%%%%%%%%%%%%

Let $\cZ = \R^{4 \times 32 \times 32}$ denote the latent space of a pretrained VAE
(e.g., the SD-VAE with scale factor $\sigma_{\text{vae}} = 0.18215$).
Let $[S] = \{1,\dots,S\}$ index style domains.

\begin{definition}[Empirical latent measures]
For a batch $B$, let $\{z_c^{(i)}\}_{i=1}^B$ be content latents and
$\{z_s^{(j)}\}_{j=1}^B$ be target-style latents.
Define the empirical measures
\begin{align}
    \hat\mu_c^B &= \frac{1}{B}\sum_{i=1}^B \delta_{z_c^{(i)}},
    &
    \hat\mu_s^B &= \frac{1}{B}\sum_{j=1}^B \delta_{z_s^{(j)}}.
\end{align}
When clear from context, we drop the superscript $B$.
\end{definition}

\begin{definition}[Transport cost]
The transport cost between two latents is the sliced-Wasserstein distance
over multi-scale patches:
\begin{equation}
    c(z, z') = \SWD_{\mathcal{R}, M}(z, z')
    = \sum_{r \in \mathcal{R}} \frac{1}{M} \sum_{m=1}^M
    \bigl\| \sort(\mathcal{P}_r(z)\,\theta_m) - \sort(\mathcal{P}_r(z')\,\theta_m) \bigr\|_2^2,
\end{equation}
where $\mathcal{P}_r(z)$ extracts patches of size $r$, $\{\theta_m\}_{m=1}^M$ are
random projection directions on the unit sphere, and $\sort$ sorts the projected
values along the spatial dimension.
In the implementation, \texttt{SWDTransportCost} wraps this computation with
micro/macro frequency decomposition and CDF-based or sort-based distance.
\end{definition}

\subsection{Code-to-Math Correspondence}

\begin{tabular}{|l|l|}
\hline
\textbf{Code} & \textbf{Math} \\
\hline
\texttt{SWDTransportCost.pairwise\_cost} & $c(z_c^{(i)}, z_s^{(j)})$ cost matrix $C_{ij}$ \\
\texttt{\_sinkhorn\_plan} & Entropic OT plan $\pi^\star = \argmin_\pi \sum_{ij}C_{ij}\pi_{ij} + \varepsilon H(\pi)$ \\
\texttt{\_sample\_from\_plan} & Matched target $\tilde z_1 \sim \pi^\star(z_s \mid z_c)$ \\
\texttt{\_bridge\_state\_and\_velocity} & Bridge state $z_t$ and target velocity $u_t$ \\
\texttt{TimeConditionedLANCETBridge.forward} & Velocity field $v_\theta(z_t, t, s)$ \\
\texttt{TimeConditionedLANCETBridge.integrate} & Euler integration $z_{k+1} = z_k + \Delta t \cdot v_\theta(z_k, t_k, s)$ \\
\texttt{w\_kinetic} & Kinetic weight $\lambda_{\text{kin}}$ \\
\texttt{terminal\_swd\_weight} & Terminal SWD weight $\lambda_{\text{term}}$ \\
\hline
\end{tabular}


%%%%%%%%%%%%%%%%%%%%%%%%%%%%%%%%%%%%%%%%%%%%%%%%%%%%%%%%%%%%%%%%%%%%%%%%%%%%%%%%
\section{Proposition 0: Empirical Latent Transport Objective}
%%%%%%%%%%%%%%%%%%%%%%%%%%%%%%%%%%%%%%%%%%%%%%%%%%%%%%%%%%%%%%%%%%%%%%%%%%%%%%%%

This proposition formalizes the complete training objective in terms of
empirical measures and latent transport, establishing the foundation for all
subsequent propositions.

\begin{proposition}[Empirical latent transport objective]
\label{prop:objective}
For a batch $B$ with content latents $\{z_c^{(i)}\}$, target-style latents
$\{z_s^{(j)}\}$, and style IDs $\{s^{(i)}\}$, the training objective is
\begin{equation}
    \cL(\theta) = \cL_{\mathrm{FM}}(\theta) + \lambda_{\text{term}} \cL_{\text{term}}(\theta) + \lambda_{\text{kin}} \cL_{\text{kin}}(\theta),
    \label{eq:full_objective}
\end{equation}
where the three components are:

\paragraph{1. Flow matching loss.}
\begin{equation}
    \cL_{\mathrm{FM}}(\theta) = \E_{z_0 \sim \hat\mu_c,\; \tilde z_1 \sim \pi^\star(\cdot \mid z_0),\; t \sim \mathcal{U}(0,1),\; \varepsilon \sim \cN(0,I)}
    \bigl\| v_\theta(z_t, t, s) - u_t \bigr\|_2^2,
\end{equation}
with bridge state and target velocity defined as
\begin{align}
    z_t &= (1-t) z_0 + t \tilde z_1 + \sigma \sqrt{t(1-t)}\,\varepsilon,
    \label{eq:bridge_state} \\
    u_t &= \tilde z_1 - z_0 + \sigma \frac{1-2t}{2\sqrt{t(1-t)}}\,\varepsilon.
    \label{eq:target_velocity}
\end{align}
Here $\pi^\star$ is the entropic OT plan solving
\begin{equation}
    \pi^\star = \argmin_{\pi \in \Pi(\hat\mu_c, \hat\mu_s)} \sum_{i,j} C_{ij}\pi_{ij} + \varepsilon H(\pi),
    \qquad C_{ij} = c(z_c^{(i)}, z_s^{(j)}),
\end{equation}
and $s$ indexes the target style domain.

\paragraph{2. Terminal SWD loss.}
Let $\Phi_\theta(z_0, s)$ denote the endpoint obtained by integrating $v_\theta$
from $z_0$ with style $s$ via Euler integration (see Proposition~3).
\begin{equation}
    \cL_{\text{term}}(\theta) = \SWD_{\mathcal{R}, M}\bigl( \Phi_\theta(z_0, s),\; \tilde z_1 \bigr).
    \label{eq:terminal_swd}
\end{equation}
In practice, $\Phi_\theta$ is computed with $N_{\text{term}} = 4$ Euler steps
during training and $N_{\text{inf}} = 12$ steps at inference.

\paragraph{3. Kinetic regularization.}
\begin{equation}
    \cL_{\text{kin}}(\theta) = \E_{z_0,\tilde z_1,t,\varepsilon}\; \|v_\theta(z_t, t, s)\|_2^2.
    \label{eq:kinetic}
\end{equation}
\end{proposition}

\begin{proof}[Derivation from code]
The objective reconstructs the \texttt{OTFlowMatchingObjective.compute()} method
(lines~727--796 of \texttt{losses.py}). The bridge state and target velocity
implement Eqs.~\eqref{eq:bridge_state}--\eqref{eq:target_velocity} exactly
(\texttt{\_bridge\_state\_and\_velocity}, lines~226--245). The OT coupling
uses \texttt{\_sinkhorn\_plan} with $L_2$-normalized SWD cost
(\texttt{SWDTransportCost.pairwise\_cost}). The terminal SWD calls
\texttt{\_terminal\_swd} which integrates $N_{\text{term}}$ steps
and computes \texttt{transport\_cost.aligned\_cost}.
The kinetic term $\cL_{\text{kin}}$ is implemented as the mean squared
velocity magnitude in \texttt{\_compute\_omf\_details} (lines~537--557).
The combined loss in the paper (Eq.~5, $\lambda_{\text{term}}=20,\lambda_{\text{kin}}=1$)
reflects this three-part structure.
\end{proof}

\begin{remark}[Two-mode implementation]
The current codebase supports two objective modes: \texttt{flow\_matching}
(time-sampled bridge states, Eq.~\eqref{eq:bridge_state}) and \texttt{omf}
(single-step endpoint prediction with $t=1$). The kinetic term appears only
in the \texttt{omf} path. The theoretical development here covers both modes
under a unified formalism: the \texttt{flow\_matching} mode provides the
time-conditioned velocity field interpretation (Propositions~1--3), while
the \texttt{omf} mode provides the complete three-part loss
(Propositions~2,~4). A future unified training mode combining all three
terms with time-sampled bridge states would realize the full objective
\eqref{eq:full_objective}.
\end{remark}


%%%%%%%%%%%%%%%%%%%%%%%%%%%%%%%%%%%%%%%%%%%%%%%%%%%%%%%%%%%%%%%%%%%%%%%%%%%%%%%%
\section{Proposition 1: Flow Matching as Conditional Velocity Estimation}
%%%%%%%%%%%%%%%%%%%%%%%%%%%%%%%%%%%%%%%%%%%%%%%%%%%%%%%%%%%%%%%%%%%%%%%%%%%%%%%%

This proposition establishes that the flow-matching objective learns the
optimal conditional expectation of the bridge velocity, connecting the
learned $v_\theta$ to the true transport dynamics.

\subsection{Theoretical statement}

\begin{assumption}[Regularity]
\label{asm:regularity}
The velocity field $v_\theta$ is parameterized by a family of measurable
functions $\cV = \{v : \cZ \times [0,1] \times [S] \to \cZ\}$ with
$\E\|v(z_t, t, s)\|^2 < \infty$ for all $(t,s)$. The bridge noise
$\varepsilon$ is independent of $(z_0, \tilde z_1, t)$.
\end{assumption}

\begin{proposition}[Bayes optimal velocity field]
\label{prop:bayes_optimal}
Under Assumption~\ref{asm:regularity}, minimizing the flow-matching objective
\begin{equation}
    \cL_{\mathrm{FM}}(\theta) = \E\bigl\| v_\theta(z_t, t, s) - u_t \bigr\|^2
\end{equation}
over $\theta$ yields the Bayes optimal solution
\begin{equation}
    v^\star(z, t, s) = \E\bigl[\, u_t \mid z_t = z,\; t,\; s \,\bigr].
    \label{eq:bayes_optimal}
\end{equation}
Furthermore, this conditional expectation decomposes as
\begin{equation}
    v^\star(z, t, s) = \E\bigl[\, \tilde z_1 - z_0 \mid z_t = z,\; t,\; s \,\bigr]
    + \sigma \frac{1-2t}{2\sqrt{t(1-t)}} \,
    \E\bigl[\, \varepsilon \mid z_t = z,\; t,\; s \,\bigr].
    \label{eq:bayes_decomp}
\end{equation}
\end{proposition}

\begin{proof}
The flow-matching objective is an $L^2$ regression problem:
$\min_\theta \E\|v_\theta(z_t, t, s) - u_t\|^2$.
For any square-integrable random vector $(X,Y)$, the minimizer of
$\E\|f(X) - Y\|^2$ over measurable $f$ is $f^\star(x) = \E[Y \mid X=x]$
(see, e.g., the standard conditional expectation as $L^2$ projection).
Applying this with $X = (z_t, t, s)$ and $Y = u_t$ gives
$v^\star(z, t, s) = \E[u_t \mid z_t = z, t, s]$.
The decomposition follows from substituting Eq.~\eqref{eq:target_velocity}
and using linearity of conditional expectation.
\end{proof}

\subsection{Implications}

\begin{corollary}[Interpretation as denoising]
\label{cor:denoising}
In the zero-noise limit $\sigma \to 0$, the bridge state degenerates to
the linear interpolation $z_t = (1-t)z_0 + t\tilde z_1$, and the optimal
velocity simplifies to
\begin{equation}
    v^\star(z, t, s) = \E\bigl[\, \tilde z_1 - z_0 \mid z_t = z,\; t,\; s \,\bigr].
\end{equation}
This is the expected transport direction conditioned on the current
interpolation state, matching the conditional flow matching formulation
of \cite{lipman2022flow}.
\end{corollary}

\begin{corollary}[Consistency with implementation signature]
\label{cor:implementation}
The learned network $v_\theta(z_t, t, s)$ approximates the conditional
expectation $\E[u_t \mid z_t, t, s]$, which depends on:
\begin{itemize}[nosep]
    \item $z_t$ (the current latent, input \texttt{x}),
    \item $t$ (the bridge time, input \texttt{t}),
    \item $s$ (the target style, input \texttt{style\_id}).
\end{itemize}
This matches the \texttt{TimeConditionedLANCETBridge.forward()} signature:
\texttt{forward(self, x, t=None, style\_id=None)}.
The time embedding is injected via sinusoidal encoding + MLP into the
style code, consistent with $t$ as a conditioning variable (see
\texttt{\_compute\_style\_code}, lines~128--149 of \texttt{model.py}).
\end{corollary}

\subsection{Limitation and bounded claim}

\begin{remark}[Approximation gap]
Proposition~\ref{prop:bayes_optimal} describes the minimizer of the
population loss under unlimited capacity. In practice:
\begin{itemize}
    \item The network has finite capacity (3.9M parameters).
    \item Training uses finite batches and stochastic optimization.
    \item The objective is optimized jointly with additional loss terms
    ($\cL_{\text{term}}, \cL_{\text{kin}}$), so the learned $v_\theta$
    is a trade-off solution rather than the pure Bayes optimal.
\end{itemize}
Therefore, the proposition provides an \emph{interpretation} of what
$v_\theta$ approximates, not a guarantee of optimality.
\end{remark}


%%%%%%%%%%%%%%%%%%%%%%%%%%%%%%%%%%%%%%%%%%%%%%%%%%%%%%%%%%%%%%%%%%%%%%%%%%%%%%%%
\section{Proposition 2: Kinetic Regularization as Path-Energy Surrogate}
%%%%%%%%%%%%%%%%%%%%%%%%%%%%%%%%%%%%%%%%%%%%%%%%%%%%%%%%%%%%%%%%%%%%%%%%%%%%%%%%

This is the most valuable theoretical connection: it establishes that
the simple squared-velocity penalty $ \cL_{\text{kin}}$ is not an ad-hoc
regularizer but a discrete proxy for the continuous-time path energy
(action functional) of the latent transport.

\subsection{Continuous-time action}

\begin{definition}[Path energy / action functional]
\label{def:action}
For a velocity field $v : \cZ \times [0,1] \times [S] \to \cZ$, define the
continuous-time action (also called the kinetic energy along paths) as
\begin{equation}
    \cA(v) = \E_{z_0, \tilde z_1, s} \left[
        \int_0^1 \bigl\| v(z_t, t, s) \bigr\|_2^2 \; dt
    \right],
    \label{eq:continuous_action}
\end{equation}
where $z_t$ follows the bridge process defined in Eq.~\eqref{eq:bridge_state}
given endpoint coupling $\pi^\star(z_0, \tilde z_1 \mid s)$.
\end{definition}

This action functional is directly analogous to the kinetic energy term
in the Benamou--Brenier formulation of optimal transport:
\begin{equation}
    \text{Benamou--Brenier:} \quad
    \inf_{(\rho, v)} \int_0^1 \int \|v_t(x)\|^2 \rho_t(x) \, dx \, dt,
\end{equation}
where $\rho_t$ is the density along the transport path. Our $\cA(v)$ is
the stochastic bridge analogue, with expectation taken over the data
distribution.

\begin{lemma}[Relationship between $\cL_{\text{kin}}$ and $\cA$]
\label{lem:kinetic_to_action}
Under Assumption~\ref{asm:regularity},
the kinetic regularizer $\cL_{\text{kin}}(\theta)$ is an unbiased estimator
of the instantaneous kinetic energy at a random time:
\begin{equation}
    \cL_{\text{kin}}(\theta) = \E_{t \sim \mathcal{U}(0,1)}\,
        \E\bigl[ \|v_\theta(z_t, t, s)\|^2 \mid t \bigr]
    = \int_0^1 \E\bigl[ \|v_\theta(z_t, t, s)\|^2 \bigr] \, dt.
    \label{eq:kinetic_as_integral}
\end{equation}
If the integrand $\E\|v_\theta(z_t, t, s)\|^2$ is constant in $t$,
then $\cL_{\text{kin}}(\theta) = \cA(v_\theta)$ exactly.
\end{lemma}

\begin{proof}
By the law of total expectation,
\begin{align}
    \cL_{\text{kin}}(\theta)
    &= \E_{z_0,\tilde z_1,t,\varepsilon} \|v_\theta(z_t, t, s)\|^2 \\
    &= \E_{t \sim \mathcal{U}(0,1)} \E_{z_0,\tilde z_1,\varepsilon}
        \bigl[ \|v_\theta(z_t, t, s)\|^2 \mid t \bigr] \\
    &= \int_0^1 \E\bigl[ \|v_\theta(z_t, t, s)\|^2 \bigr] \, dt \quad
        \text{(since $t \sim \mathcal{U}(0,1)$)}.
\end{align}
Under the assumption that $\E\|v_\theta(z_t, t, s)\|^2$ is independent of $t$,
this integral equals $\cA(v_\theta)$ by Definition~\ref{def:action}.
\end{proof}

\begin{remark}[When $t$-independence holds]
The assumption of constant $\E\|v_\theta\|^2$ across $t$ is approximately
satisfied when:
\begin{itemize}
    \item The velocity field is learned to produce roughly straight paths
    (as encouraged by the flow matching objective, which targets $u_t$ with
    variance that depends on $t$ but mean $\tilde z_1 - z_0$).
    \item The training samples $t$ uniformly, so the network cannot specialize
    its velocity magnitude to specific time intervals.
\end{itemize}
Empirically, the near-flat step-count sweep (4--16 steps give similar quality)
suggests that the learned dynamics are approximately uniform, consistent with
this assumption.
\end{remark}

\subsection{Discrete path energy at inference}

\begin{proposition}[Discrete action approximation]
\label{prop:discrete_action}
Let $\{z_k\}_{k=0}^K$ be the Euler integration trajectory (see
Proposition~\ref{prop:euler_error}) with step size $\Delta t = 1/K$.
Define the discrete path energy
\begin{equation}
    \cA_{\Delta t}(v_\theta) = \sum_{k=0}^{K-1} \Delta t \,
        \bigl\| v_\theta(z_k, t_k, s) \bigr\|_2^2,
    \qquad t_k = (k + 0.5)\, \Delta t.
    \label{eq:discrete_action}
\end{equation}
Under the regularity conditions of Proposition~\ref{prop:euler_error}
($v_\theta$ is $L$-Lipschitz in $z$ and bounded by $M$), we have the
pointwise approximation bound for each trajectory:
\begin{equation}
    \bigl| \cA_{\Delta t}(v_\theta) - \cA(v_\theta) \bigr| \leq C \Delta t,
    \label{eq:discrete_action_bound}
\end{equation}
where $C$ depends on $L$, $M$, and the trajectory length.
\end{proposition}

\begin{proof}[Sketch]
The continuous action $\cA(v_\theta)$ is the integral
$\int_0^1 \|v_\theta(z(t), t, s)\|^2 \, dt$.
The discrete action $\cA_{\Delta t}$ is the Riemann sum approximation
with midpoint rule. The error between the true trajectory $z(t)$ and the
Euler trajectory $z_k$ accumulates at $O(\Delta t)$ (Proposition~3), and
since $v_\theta$ is Lipschitz, the integrand error is also $O(\Delta t)$.
Standard numerical analysis gives $O(\Delta t)$ convergence of the midpoint
Riemann sum for Lipschitz integrands.
\end{proof}

\subsection{Code correspondence}

\begin{itemize}
    \item The kinetic loss is computed as \texttt{pred\_velocity.float()\ **\ 2}
    followed by \texttt{.mean()} (lines~537--548 of \texttt{losses.py}).
    \item The weight \texttt{w\_kinetic = 1.0} corresponds to
    $\lambda_{\text{kin}} = 1$ in the paper, meaning $\cL_{\text{kin}}$
    enters the total loss at unit scale.
    \item An optional entropy-gated variant (\texttt{kinetic\_entropy\_gate\_weight})
    spatially modulates the kinetic penalty by the semantic attention entropy,
    but the mainline uses the uniform version described here.
\end{itemize}

\begin{remark}[Conceptual value]
Proposition~\ref{prop:discrete_action} transforms the kinetic penalty from
a heuristic magnitude constraint into a principled path-energy regularizer.
This connects our training to the dynamical optimal transport literature,
where the action $\int\|v\|^2$ controls the transport cost and path regularity.
The discrete approximation bound justifies why the penalty is applied
per-step rather than on the integrated endpoint.
\end{remark}


%%%%%%%%%%%%%%%%%%%%%%%%%%%%%%%%%%%%%%%%%%%%%%%%%%%%%%%%%%%%%%%%%%%%%%%%%%%%%%%%
\section{Proposition 3: Euler Integration Discretization Error}
%%%%%%%%%%%%%%%%%%%%%%%%%%%%%%%%%%%%%%%%%%%%%%%%%%%%%%%%%%%%%%%%%%%%%%%%%%%%%%%%

This proposition provides a standard numerical analysis bound for the
Euler integration used at inference, directly supporting the step-count
experiments in the paper.

\subsection{Problem setup}

Consider the learned velocity field $v_\theta : \cZ \times [0,1] \times [S] \to \cZ$.
At inference, we solve the initial value problem (IVP):
\begin{equation}
    \frac{d}{dt} z(t) = v_\theta(z(t), t, s), \qquad z(0) = z_0,
    \label{eq:ivp}
\end{equation}
using explicit Euler integration with $K$ steps:
\begin{equation}
    z_0^{\text{(disc)}} = z_0, \qquad
    z_{k+1}^{\text{(disc)}} = z_k^{\text{(disc)}} + \Delta t \,
        v_\theta\bigl(z_k^{\text{(disc)}}, t_k, s\bigr),
    \qquad \Delta t = \frac{1}{K},\; t_k = (k + 0.5)\, \Delta t.
    \label{eq:euler}
\end{equation}
The final point is $z_K^{\text{(disc)}} \approx z(1)$.

\begin{assumption}[Lipschitz and boundedness]
\label{asm:lipschitz}
The velocity field $v_\theta$ satisfies:
\begin{enumerate}[nosep]
    \item \textbf{Lipschitz in $z$:}
    $\|v_\theta(z_1, t, s) - v_\theta(z_2, t, s)\| \leq L \|z_1 - z_2\|$
    for all $z_1, z_2 \in \cZ$, $t \in [0,1]$, $s \in [S]$.
    \item \textbf{Uniform boundedness:}
    $\sup_{z,t,s} \|v_\theta(z, t, s)\| \leq M < \infty$.
    \item \textbf{Lipschitz in $t$:}
    $\|v_\theta(z, t_1, s) - v_\theta(z, t_2, s)\| \leq L_t |t_1 - t_2|$
    for all $z \in \cZ$, $t_1, t_2 \in [0,1]$, $s \in [S]$.
\end{enumerate}
\end{assumption}

\begin{proposition}[Global truncation error]
\label{prop:euler_error}
Under Assumption~\ref{asm:lipschitz}, the Euler integration error at the
terminal time $t=1$ satisfies
\begin{equation}
    \bigl\| z(1) - z_K^{\text{(disc)}} \bigr\|
    \leq \frac{\Delta t}{2} \bigl( M L_t + M^2 L \bigr)
        \frac{e^L - 1}{L}
    \leq C \cdot \Delta t,
    \label{eq:euler_error_bound}
\end{equation}
where $C = \frac{1}{2} (M L_t + M^2 L) \frac{e^L - 1}{L}$ depends on the
velocity field regularity but not on $K$.
\end{proposition}

\begin{proof}
We follow the standard global truncation error analysis for Euler's method
on ODEs (see, e.g., \cite{atkinson2009numerical}).

\paragraph{Local truncation error.}
Let $z(t)$ be the true solution and $z_k$ the discrete approximation.
The local truncation error (LTE) at step $k$ is:
\begin{align}
    \tau_k &= z(t_{k+1}) - \bigl[ z(t_k) + \Delta t \, v_\theta(z(t_k), t_k, s) \bigr] \\
    &= \int_{t_k}^{t_{k+1}} v_\theta(z(\tau), \tau, s) \, d\tau
        - \Delta t \, v_\theta(z(t_k), t_k, s) .
\end{align}
Applying Taylor expansion with integral remainder:
\begin{align}
    \|\tau_k\| &\leq \int_{t_k}^{t_{k+1}} \bigl\| v_\theta(z(\tau), \tau, s) - v_\theta(z(t_k), t_k, s) \bigr\| \, d\tau \\
    &\leq \int_{t_k}^{t_{k+1}} \bigl( L \|z(\tau) - z(t_k)\| + L_t |\tau - t_k| \bigr) \, d\tau \\
    &\leq \int_{t_k}^{t_{k+1}} \bigl( L M |\tau - t_k| + L_t |\tau - t_k| \bigr) \, d\tau \\
    &= (L M + L_t) \int_0^{\Delta t} s \, ds = \frac{1}{2} (L M + L_t) (\Delta t)^2 .
\end{align}
Thus $\|\tau_k\| \leq \frac{1}{2} (L M + L_t) (\Delta t)^2$.

\paragraph{Global error propagation.}
Let $e_k = \|z(t_k) - z_k\|$. Then:
\begin{align}
    e_{k+1} &\leq \|z(t_{k+1}) - [z_k + \Delta t \, v_\theta(z_k, t_k, s)]\| \\
    &\leq \|z(t_k) - z_k\| + \Delta t \|v_\theta(z(t_k), t_k, s) - v_\theta(z_k, t_k, s)\| + \|\tau_k\| \\
    &\leq e_k + \Delta t \, L \, e_k + \frac{1}{2} (L M + L_t) (\Delta t)^2 .
\end{align}

\paragraph{Solving the recurrence.}
We have $e_{k+1} \leq (1 + L \Delta t) e_k + \frac{1}{2} (L M + L_t) (\Delta t)^2$.
With $e_0 = 0$, this gives:
\begin{align}
    e_K &\leq \frac{1}{2} (L M + L_t) (\Delta t)^2 \sum_{j=0}^{K-1} (1 + L \Delta t)^j \\
    &\leq \frac{1}{2} (L M + L_t) (\Delta t)^2 \frac{(1 + L \Delta t)^K - 1}{L \Delta t} \\
    &\leq \frac{1}{2} (L M + L_t) \Delta t \frac{e^{L K \Delta t} - 1}{L} \\
    &= \frac{1}{2} (L M + L_t) \Delta t \frac{e^L - 1}{L} .
\end{align}
Here we used the inequality $(1 + L\Delta t)^K \leq e^{L K \Delta t}$,
and $K \Delta t = 1$. The terminal error is $e_K$, giving the bound.
\end{proof}

\subsection{Implications for step-count experiments}

\begin{corollary}[Consistency with observed flatness]
\label{cor:step_flatness}
The $O(\Delta t)$ error bound explains the near-flat quality across
4--16 steps observed in the paper (Section~5.3, Step-count sweep).
Specifically:
\begin{itemize}
    \item At $K = 4$ steps ($\Delta t = 0.25$), the bound gives
    $e_K \leq C / 4$.
    \item At $K = 16$ steps ($\Delta t = 0.0625$), the bound gives
    $e_K \leq C / 16$.
    \item If $C$ is small (well-conditioned velocity field), the absolute
    error at 4 steps is already small, so further refinement yields
    diminishing returns.
\end{itemize}
The reported CLIP-style variation of only 0.7159--0.7162 and LPIPS
variation of $\sim 0.4514$ across $K \in \{1, 4, 8, 12, 16\}$ is
consistent with a well-conditioned $v_\theta$ in the trained operating
regime.
\end{corollary}

\begin{remark}[Constants are data-dependent]
The constants $L$, $M$, and $L_t$ in Assumption~\ref{asm:lipschitz} depend
on the learned parameters $\theta$ and the latent data distribution.
Proposition~\ref{prop:euler_error} does not provide a quantitative bound
applicable to any specific checkpoint; rather, it establishes that the
error scales as $O(\Delta t)$, with the constant determined by the
regularity of the learned field. The empirical flatness across steps
indicates that the operating $C$ is sufficiently small.
\end{remark}


%%%%%%%%%%%%%%%%%%%%%%%%%%%%%%%%%%%%%%%%%%%%%%%%%%%%%%%%%%%%%%%%%%%%%%%%%%%%%%%%
\section{Proposition 4: Terminal SWD as Endpoint Distribution Control}
%%%%%%%%%%%%%%%%%%%%%%%%%%%%%%%%%%%%%%%%%%%%%%%%%%%%%%%%%%%%%%%%%%%%%%%%%%%%%%%%

This proposition formalizes the terminal SWD loss as a mechanism for
controlling the distribution of generated endpoint patches, rather than
matching individual target images.

\subsection{Preliminaries: Patch distributions}

\begin{definition}[Patch empirical measure]
For a latent $z \in \cZ$ and patch size $r \in \mathbb{N}$, let
$\mathcal{P}_r(z) \subset \R^c$ be the set of all $c$-channel patches
of spatial size $r \times r$ extracted from $z$ (with stride 1).
For a batch of generated endpoints $\{\hat z_1^{(i)}\}$ and target-style
latents $\{z_s^{(j)}\}$, define the empirical patch measures:
\begin{align}
    \hat\nu_\theta^{(r)} &= \frac{1}{N_\theta} \sum_{i=1}^B \sum_{p \in \mathcal{P}_r(\hat z_1^{(i)})} \delta_p,
    &
    \hat\nu_s^{(r)} &= \frac{1}{N_s} \sum_{j=1}^B \sum_{p \in \mathcal{P}_r(z_s^{(j)})} \delta_p,
\end{align}
where $N_\theta$ and $N_s$ are the total number of patches extracted.
\end{definition}

\begin{definition}[Sliced Wasserstein Distance]
\label{def:swd}
For two probability measures $\mu, \nu$ on $\R^c$ with finite first moments,
the sliced Wasserstein distance with $M$ projections is
\begin{equation}
    \SWD_M(\mu, \nu) = \frac{1}{M} \sum_{m=1}^M
        \bigl\| \sort(\langle \mu, \theta_m\rangle) - \sort(\langle \nu, \theta_m\rangle) \bigr\|_2,
    \label{eq:swd_definition}
\end{equation}
where $\{\theta_m\}_{m=1}^M$ are uniformly sampled from the unit sphere
$\mathbb{S}^{c-1}$, and $\sort(\langle \mu, \theta\rangle)$ denotes the
sorted list of one-dimensional projections of samples from $\mu$ onto
$\theta$.
\end{definition}

\subsection{Theoretical statement}

\begin{proposition}[Terminal SWD controls patch distribution]
\label{prop:terminal_swd}
The terminal SWD loss in the paper,
\begin{equation}
    \cL_{\text{term}}(\theta) = \frac{1}{|\mathcal{R}|} \sum_{r \in \mathcal{R}}
        \SWD_M\bigl( \hat\nu_\theta^{(r)},\; \hat\nu_s^{(r)} \bigr),
    \qquad \mathcal{R} = \{3, 5, 7, 15\},
    \label{eq:terminal_swd_loss}
\end{equation}
satisfies the following properties:

\begin{enumerate}[nosep]
    \item \textbf{Identifiability:}
    If $\cL_{\text{term}}(\theta) = 0$, then for every $r \in \mathcal{R}$
    and every projection direction $\theta_m$, the empirical sorted
    projections of $\hat\nu_\theta^{(r)}$ and $\hat\nu_s^{(r)}$ are identical.
    In the limit $M \to \infty$, this implies
    $\hat\nu_\theta^{(r)} \overset{d}{=} \hat\nu_s^{(r)}$ for all $r \in \mathcal{R}$.

    \item \textbf{Content-style disentanglement:}
    The loss operates on patch statistics rather than spatial correspondence.
    Two latents $z_1, z_2$ can have identical patch distributions
    (i.e., $\SWD_M(\mathcal{P}_r(z_1), \mathcal{P}_r(z_2)) = 0$)
    while having very different spatial arrangements of patches.
    This is precisely the desired behavior for style transfer: the
    texture statistics should match the target style, but the content
    structure should be preserved from the source.

    \item \textbf{Multi-scale control:}
    By averaging over patch sizes $\{3,5,7,15\}$, the loss controls
    statistics at multiple scales: small patches ($r=3$) capture fine
    texture, while larger patches ($r=15$) capture tonal distribution
    and coarse pattern statistics.
\end{enumerate}
\end{proposition}

\begin{proof}
\paragraph{Property 1.}
If $\SWD_M(\hat\nu_\theta^{(r)}, \hat\nu_s^{(r)}) = 0$,
then for each projection $m$, $\|\sort(\langle\hat\nu_\theta^{(r)},\theta_m\rangle) - \sort(\langle\hat\nu_s^{(r)},\theta_m\rangle)\| = 0$.
Thus the sorted one-dimensional projections are identical for all $M$
directions. As $M \to \infty$, the SWD converges to the Wasserstein-1
distance on the sliced Radon transform \cite{bonnotte2013non}, and zero
Wasserstein-1 distance implies distributional equality.

\paragraph{Property 2.}
Patch extraction is permutation-invariant: $\mathcal{P}_r(z)$ is a
multiset of patches without spatial indexing. Therefore, $\SWD$
measures only the distribution of patch features, not their spatial
layout. Two images with identical patch bags (e.g., a photo and its
patch-shuffled version) have $\SWD = 0$ despite different spatial
structure. This is precisely the style-content separation property
exploited in texture synthesis \cite{julesz1981textons}.

\paragraph{Property 3.}
The patch size $r$ controls the receptive field of the SWD comparison.
Small $r$ captures high-frequency texture details; large $r$ captures
low-frequency tonal layout. The multi-scale average in Eq.~\eqref{eq:terminal_swd_loss}
ensures the endpoint matches the target distribution at all relevant
texture scales.
\end{proof}

\begin{corollary}[Ablation interpretation]
\label{cor:ablation_terminal}
The destructive ablation ``D1: w/o terminal SWD'' (paper Table~3) removes
$\cL_{\text{term}}$, which by Proposition~\ref{prop:terminal_swd} removes
the mechanism for controlling the endpoint patch distribution toward the
target domain. The observed degradation in CLIP-style (0.6708 vs.\ 0.7014)
and improvement in LPIPS (0.3490 vs.\ 0.4593) is consistent with the
interpretation that without terminal SWD, the model defaults to content
preservation (low LPIPS) but fails to acquire target texture statistics
(low CLIP-style).
\end{corollary}

\begin{remark}[Bounded claim]
Proposition~\ref{prop:terminal_swd} does not claim that $\cL_{\text{term}}=0$
implies the full generated endpoint distribution matches the full target
distribution in all respects. The SWD is:
\begin{itemize}
    \item Limited to $M = 64$ projections (far short of the $M\to\infty$ limit).
    \item Computed on instance-normalized features, not raw latents.
    \item Restricted to patch sizes $\{3,5,7,15\}$, missing both
    single-pixel ($r=1$) and global ($r=32$) statistics.
\end{itemize}
The loss is a \emph{practical proxy} for distribution matching, not an
exact transport solver. This bounded interpretation is stated explicitly
in the paper (Section~3.3, ``Bounded theoretical claim'').
\end{remark}


%%%%%%%%%%%%%%%%%%%%%%%%%%%%%%%%%%%%%%%%%%%%%%%%%%%%%%%%%%%%%%%%%%%%%%%%%%%%%%%%
\section{Proposition 5: OT Coupling Stabilizes Endpoint Supervision}
%%%%%%%%%%%%%%%%%%%%%%%%%%%%%%%%%%%%%%%%%%%%%%%%%%%%%%%%%%%%%%%%%%%%%%%%%%%%%%%%

This is the most method-specific proposition, connecting the OT coupling
to the stability of the learned transport.

\subsection{Motivation}

In the absence of paired data, the flow matching objective requires a
supervision endpoint $\tilde z_1$ for each content latent $z_0$.
Two coupling strategies are available:
\begin{enumerate}[nosep]
    \item \textbf{Random coupling:} uniformly sample $\tilde z_1$ from
    the target-style pool.
    \item \textbf{OT coupling:} sample $\tilde z_1 \sim \pi^\star(z_s \mid z_c)$
    minimizing the SWD transport cost.
\end{enumerate}
We analyze how the coupling choice affects the supervision signal.

\subsection{Theoretical development}

\begin{definition}[Coupling cost and supervision variance]
For a coupling $\pi$ over a batch $(Z_c, Z_s) \sim (\hat\mu_c, \hat\mu_s)$,
define:
\begin{align}
    \mathcal{C}(\pi) &= \E_{(z_0, \tilde z_1) \sim \pi} \bigl[ c(z_0, \tilde z_1) \bigr]
        \quad \text{(expected transport cost)}, \\
    \mathcal{V}(\pi) &= \Var_{(z_0, \tilde z_1) \sim \pi} \bigl[ \tilde z_1 - z_0 \bigr]
        \quad \text{(supervision variance)}.
\end{align}
\end{definition}

\begin{proposition}[OT coupling minimizes cost]
\label{prop:ot_min_cost}
For a fixed batch, let $\Pi(\hat\mu_c, \hat\mu_s)$ be the set of all
couplings between the empirical measures. The entropic OT coupling
\begin{equation}
    \pi^\star = \argmin_{\pi \in \Pi} \sum_{i,j} C_{ij} \pi_{ij} + \varepsilon H(\pi)
\end{equation}
satisfies $\mathcal{C}(\pi^\star) \leq \mathcal{C}(\pi)$ for any
deterministic coupling $\pi$ (up to the entropy bias $\varepsilon H$).
In the limit $\varepsilon \to 0$, $\pi^\star$ recovers the exact OT coupling
with minimal expected cost.
\end{proposition}

\begin{proof}
By definition, $\pi^\star$ minimizes $\sum C_{ij}\pi_{ij} + \varepsilon H(\pi)$.
For any deterministic coupling $\pi$ (a permutation matrix), $H(\pi)=0$, so
$\mathcal{C}(\pi^\star) + \varepsilon H(\pi^\star) \leq \mathcal{C}(\pi)$.
Thus $\mathcal{C}(\pi^\star) \leq \mathcal{C}(\pi)$ since $H(\pi^\star) \geq 0$.
As $\varepsilon \to 0$, the entropic regularization vanishes and $\pi^\star$
converges to a deterministic OT plan (under non-degenerate cost), giving
$\mathcal{C}(\pi^\star) \leq \mathcal{C}(\pi)$ for all $\pi$.
\end{proof}

\begin{proposition}[Lower cost implies lower velocity variance]
\label{prop:cost_to_variance}
Assume a simplified setting where $v_\theta$ is linear in the endpoints:
$v_\theta(z_t, t, s) \approx \E[\tilde z_1 - z_0 \mid z_t, t, s]$
(ignoring the noise term in $u_t$ for analytical tractability).
Under a linear bridge model with $z_t = (1-t)z_0 + t\tilde z_1$,
the conditional target velocity given $z_t$ has variance bounded by
the coupling cost:
\begin{equation}
    \Var\bigl[ \tilde z_1 - z_0 \mid z_t, t, s \bigr]
    \leq \frac{1}{t^2} \Var\bigl[ \tilde z_1 \mid z_0 \bigr].
\end{equation}
\end{proposition}

\begin{proof}[Sketch]
In the linear bridge ($\sigma=0$), $z_t = (1-t)z_0 + t\tilde z_1$ and the
target velocity is $u_t = \tilde z_1 - z_0$. The conditional distribution
of $\tilde z_1 - z_0$ given $z_t$ is determined by the joint of $(z_0, \tilde z_1)$.
A standard Gaussian conditioning argument gives the variance bound.
\end{proof}

\begin{corollary}[OT coupling reduces supervision noise]
\label{cor:ot_stabilization}
Combining Propositions~\ref{prop:ot_min_cost} and~\ref{prop:cost_to_variance},
the OT coupling reduces both the expected cost and (under the linear model)
the target velocity variance compared to random coupling. This means:
\begin{itemize}
    \item Each training step sees more consistent target velocities.
    \item The flow matching objective has lower conditional variance,
    making optimization easier.
    \item The learned velocity field is smoother and more stable,
    consistent with the near-flat step-count curve.
\end{itemize}
\end{corollary}

\begin{remark}[Empirical verification pathway]
This proposition naturally pairs with targeted experiments:
\begin{itemize}
    \item Compare training loss curves: OT vs.\ random coupling
    ("random shuffle" baseline should show higher variance).
    \item Measure trajectory curvature/straightness under both couplings
    (OT coupling should produce straighter paths).
    \item Compare target velocity variance during training.
    \item Compare endpoint mismatch under both couplings.
\end{itemize}
Each of these comparisons would strengthen the theoretical claim and
connect it to observable quantities.
\end{remark}

\begin{remark}[Not a complete proof]
Proposition~\ref{prop:cost_to_variance} relies on a simplified linear model
($\sigma=0$) and approximates $v_\theta$ by its Bayes optimal form.
The actual bridge includes noise ($\sigma > 0$), and the learned
$v_\theta$ is a finite-capacity approximation. The corollary provides
a plausible mechanism for OT coupling's benefit rather than a rigorous
guarantee. A complete analysis would require:
\begin{itemize}
    \item A precise characterization of $\Var[\tilde z_1 - z_0 \mid z_t, t, s]$
    under the full stochastic bridge ($\sigma > 0$).
    \item A bound connecting coupling cost to optimization convergence rate.
\end{itemize}
These extensions are left for future work.
\end{remark}


%%%%%%%%%%%%%%%%%%%%%%%%%%%%%%%%%%%%%%%%%%%%%%%%%%%%%%%%%%%%%%%%%%%%%%%%%%%%%%%%
\section{Unified Discussion}
%%%%%%%%%%%%%%%%%%%%%%%%%%%%%%%%%%%%%%%%%%%%%%%%%%%%%%%%%%%%%%%%%%%%%%%%%%%%%%%%

\subsection{Summary of contributions}

The five propositions form a coherent theory package:

\begin{enumerate}[nosep]
    \item \textbf{Proposition 0} establishes the formal objective connecting
    code to math.
    \item \textbf{Proposition 1} interprets $v_\theta$ as the Bayes optimal
    conditional velocity estimator.
    \item \textbf{Proposition 2} justifies $\cL_{\text{kin}}$ as a
    path-energy surrogate with discrete approximation.
    \item \textbf{Proposition 3} bounds the Euler integration error
    at $O(\Delta t)$, explaining the step-count flatness.
    \item \textbf{Proposition 4} formalizes terminal SWD as endpoint
    distribution control.
    \item \textbf{Proposition 5} analyzes OT coupling's stabilizing
    effect on endpoint supervision.
\end{enumerate}

\subsection{What these results do and do not claim}

\begin{tabular}{|l|l|}
\hline
\textbf{The results show} & \textbf{The results do not show} \\
\hline
A formalization of the training objective & Equivalence to full Schr\"odinger Bridge \\
Conditional velocity interpretation & Global optimal transport solution \\
Kinetic term $\leftrightarrow$ path energy & Optimality of the multi-scale SWD \\
Euler error $O(\Delta t)$ scaling & Tightness of the error constants \\
Terminal SWD $\leftrightarrow$ patch distribution & Weak convergence guarantees \\
OT coupling $\leftrightarrow$ lower cost \& variance & Optimality of Sinkhorn for this task \\
\hline
\end{tabular}

\subsection{Recommended paper integration}

We recommend adding a new subsection in the Method section:

\begin{quote}
\textbf{3.X Formal analysis as empirical latent transport}
\begin{enumerate}[nosep]
    \item 3.X.1 Empirical latent transport objective (Proposition~0)
    \item 3.X.2 Flow matching as conditional velocity estimation (Proposition~1)
    \item 3.X.3 Kinetic regularization as path-energy surrogate (Proposition~2)
    \item 3.X.4 Euler discretization error of latent integration (Proposition~3)
    \item 3.X.5 Terminal SWD as endpoint distribution control (Proposition~4)
    \item 3.X.6 Why OT coupling stabilizes endpoint supervision (Proposition~5, optional)
\end{enumerate}
\end{quote}

Each sub-subsection should include:
\begin{itemize}
    \item One explicit proposition with clear assumptions.
    \item A concise proof or proof sketch (appendix for full proofs).
    \item An explicit ``Code correspondence'' paragraph connecting the math
    to specific implementation details.
    \item A ``Bounded claim'' remark noting limitations.
\end{itemize}

\subsection{Checklist implications}

Based on the above, the reproducibility checklist item
``Does this paper make theoretical contributions?'' could be changed to
\textbf{Yes} with the following justification:

\begin{quote}
\textbf{Yes (bounded).} The paper provides a formal analysis of the proposed
latent transport framework, including: the Bayes optimal interpretation of
the flow matching objective (Proposition~1), the path-energy interpretation
of the kinetic regularizer with discrete approximation bounds (Proposition~2),
the $O(\Delta t)$ Euler discretization error analysis (Proposition~3),
and the endpoint distribution control interpretation of terminal SWD
(Proposition~4). These results do not claim exact equivalence to a full
Schr\"odinger Bridge or Benamou--Brenier solver; they provide a bounded
formal justification for the three-part design.
\end{quote}

\subsection{Open questions and future directions}

\begin{enumerate}
    \item \textbf{Unified training objective:} The current code separates
    the time-sampled flow matching (\texttt{flow\_matching}) and the
    three-part loss (\texttt{omf}) into different modes. A unified training
    objective combining time-sampled bridge states with kinetic regularization
    and terminal SWD would realize the full objective of Eq.~\eqref{eq:full_objective}
    and strengthen the theoretical-to-empirical connection.

    \item \textbf{Tightness of error bounds:} The constants $L, M, L_t$ in
    Proposition~\ref{prop:euler_error} are unknown for the trained model.
    Empirically estimating these constants from the learned $v_\theta$
    would provide quantitative bounds.

    \item \textbf{Distributional convergence of terminal SWD:} The SWD
    convergence rate as $M \to \infty$ is $\Theta(M^{-1/2})$ for fixed
    sample size \cite{nietert2022wasserstein}, but the joint scaling of
    $M$ and batch size $B$ for the empirical SWD is not fully characterized.

    \item \textbf{Coupling cost and generalization:} Proposition~5 suggests
    OT coupling reduces supervision variance, but the precise relationship
    between coupling cost and generalization error in the learned transport
    remains an open question.
\end{enumerate}

\bibliographystyle{plain}
\begin{thebibliography}{10}

\bibitem{atkinson2009numerical}
K.~Atkinson, W.~Han, and D.~Stewart.
\newblock {\em Numerical Solution of Ordinary Differential Equations}.
\newblock Wiley, 2009.

\bibitem{bonnotte2013non}
N.~Bonnotte.
\newblock {\em Non-Asymptotic Bounds for the Sliced Wasserstein Distance}.
\newblock PhD thesis, Universit\'e Paris-Sud, 2013.

\bibitem{julesz1981textons}
B.~Julesz.
\newblock Textons, the elements of texture perception, and their interactions.
\newblock {\em Nature}, 290(5802):91--97, 1981.

\bibitem{lipman2022flow}
Y.~Lipman, R.~T.~Q.~Chen, H.~Ben-Hamu, M.~Nickel, and M.~Le.
\newblock Flow matching for generative modeling.
\newblock In {\em ICLR}, 2023.

\bibitem{nietert2022wasserstein}
S.~Nietert, Z.~Goldfeld, and K.~Kato.
\newblock Statistical and computational limits of sliced Wasserstein distances.
\newblock In {\em NeurIPS}, 2022.

\end{thebibliography}

\end{document}
